\def\paperroot{paper/}
\IfFileExists{preamble.tex}{\def\paperroot{}}{}

\input{\paperroot preamble}

\begin{document}

\begin{abstract}
For a finite set $U\subset\mathbb{R}^{2}$, let $u(U)$ denote the number
of ordered pairs of points in $U$ at distance exactly $1$. In this
paper we improve upon the best existing explicit lower bound for the
Erd\H{o}s unit distance problem, due to Sawin \cite{Sawin}, and prove
that there are arbitrarily large finite sets $U\subset\mathbb{R}^{2}$
satisfying
\[
u(U)\geq C^{-1}(\#U)^{1.0358324},
\]
where $C>0$ is an absolute constant. The construction follows the
number field strategy introduced in the recent disproof of Erd\H{o}s's
unit distance conjecture, where unit distances arise from elements of
prescribed relative norm in an ideal of a CM field, and the fields are
supplied by a Golod-Shafarevich tower over a real quadratic field
$\mathbb{Q}(\sqrt{D})$. The main new ingredient is a tower that is
allowed to ramify at the unique dyadic prime of $\mathbb{Q}(\sqrt{D})$,
which makes primes that split in $\mathbb{Q}(\sqrt{D})$ affordable in
the Golod-Shafarevich criterion.
\end{abstract}

\maketitle

% AI methodology statement, referenced by the asterisk on the author's name.
{\noindent\small\textsuperscript{*}\textbf{AI Methodology.} The research
for this paper was carried out with GPT 5.5 Pro, and the paper was written
with Codex 5.6 Sol, under the supervision of the author. The author has read
all of the arguments carefully.\par}
\medskip

\input{\paperroot chapters/01_introduction}
\input{\paperroot chapters/02_norm_count_pigeonhole}
\input{\paperroot chapters/03_louboutin_master}
\input{\paperroot chapters/04_dyadic_tower}
\input{\paperroot chapters/05_certificate}

\input{\paperroot appendices/selected_prime_certificate}
\input{\paperroot appendices/ramified_prime_witnesses}

\providecommand{\bysame}{\leavevmode\hbox to3em{\hrulefill}\thinspace}
\providecommand{\MR}{\relax\ifhmode\unskip\space\fi MR }
% \MRhref is called by the amsart/book/proc definition of \MR.
\providecommand{\MRhref}[2]{%
  \href{http://www.ams.org/mathscinet-getitem?mr=#1}{#2}
}
\providecommand{\href}[2]{#2}
\begin{thebibliography}{10}

\bibitem{AlonEtAl}
Noga Alon, Thomas~F. Bloom, W.~Timothy Gowers, Daniel Litt, Will Sawin, Arul
  Shankar, Jacob Tsimerman, Victor Wang, and Melanie~Matchett Wood,
  \emph{Remarks on the disproof of the unit distance conjecture}, 2026,
  arXiv:2605.20695 [math.CO].

\bibitem{Banaszczyk}
Wojciech Banaszczyk, \emph{New bounds in some transference theorems in the
  geometry of numbers}, Mathematische Annalen \textbf{296} (1993), 625--635.

\bibitem{Erdos1946}
Paul Erd{\H o}s, \emph{On sets of distances of {$n$} points}, The American
  Mathematical Monthly \textbf{53} (1946), no.~5, 248--250.

\bibitem{Erdos1994}
\bysame, \emph{Some problems in number theory, combinatorics and combinatorial
  geometry}, Mathematica Pannonica \textbf{5} (1994), no.~2, 261--269.

\bibitem{GolodShafarevich}
Evgenii~S. Golod and Igor~R. Shafarevich, \emph{On the class field tower}, Izv.
  Akad. Nauk SSSR Ser. Mat. \textbf{28} (1964), 261--272.

\bibitem{HajirMaireRamakrishna}
Farshid Hajir, Christian Maire, and Ravi Ramakrishna, \emph{On the
  {Shafarevich} group of restricted ramification extensions of number fields in
  the tame case}, Indiana University Mathematics Journal \textbf{70} (2021),
  no.~6, 2693--2710.

\bibitem{HalberstamRichert}
Heini Halberstam and Hans-Egon Richert, \emph{Sieve methods}, Academic Press,
  1974.

\bibitem{IwaniecKowalski}
Henryk Iwaniec and Emmanuel Kowalski, \emph{Analytic number theory}, AMS
  Colloquium Publications, vol.~53, American Mathematical Society, 2004.

\bibitem{Koch}
Helmut Koch, \emph{Zum {Satz} von {Golod-Schafarewitsch}}, Mathematische
  Nachrichten \textbf{42} (1969), 321--333.

\bibitem{Louboutin}
St{\'e}phane Louboutin, \emph{Explicit bounds for residues of {Dedekind} zeta
  functions, values of {$L$}-functions at {$s=1$}, and relative class numbers},
  Journal of Number Theory \textbf{85} (2000), 263--282.

\bibitem{MontgomeryVaughan}
Hugh~L. Montgomery and Robert~C. Vaughan, \emph{The large sieve}, Mathematika
  \textbf{20} (1973), no.~2, 119--134.

\bibitem{NSW}
J{\"u}rgen Neukirch, Alexander Schmidt, and Kay Wingberg, \emph{Cohomology of
  number fields}, 2 ed., Grundlehren der mathematischen Wissenschaften, vol.
  323, Springer, 2008.

\bibitem{OpenAIUnitDistance}
{OpenAI}, \emph{Planar point sets with many unit distances}, 2026,
  URL:https://openai.com/news/.

\bibitem{Sawin}
Will Sawin, \emph{An explicit lower bound for the unit distance problem}, 2026,
  arXiv:2605.20579 [math.CO].

\bibitem{Schmidt}
Alexander Schmidt, \emph{On the relation between {$2$} and {$\infty$} in
  {Galois} cohomology of number fields}, Compositio Mathematica \textbf{133}
  (2002), 267--288.

\bibitem{SpencerSzemerediTrotter}
Joel Spencer, Endre Szemer{\'e}di, and William~T. Trotter, Jr., \emph{Unit
  distances in the {Euclidean} plane}, Graph Theory and Combinatorics
  (Cambridge, 1983), Academic Press, London, 1984, pp.~293--303.

\bibitem{Spiderduckpig}
{Spiderduckpig}, \emph{What is the unit distance exponent?}, 2026, MathOverflow
  answer, URL (accessed on 2026-07-06): https://mathoverflow.net/a/511531.

\bibitem{Vinberg}
Ernest~B. Vinberg, \emph{On the theorem concerning the infinite-dimensionality
  of an associative algebra}, Izv. Akad. Nauk SSSR Ser. Mat. \textbf{29}
  (1965), 209--214.

\bibitem{Washington}
Lawrence~C. Washington, \emph{Introduction to cyclotomic fields}, 2 ed.,
  Graduate Texts in Mathematics, vol.~83, Springer, 1997.

\end{thebibliography}

\end{document}
